% Detailed derivation of Eqs. (7)-(14) of theory-note.tex.
% Draft of 2026-11-20. Equation numbers of the theory note refer to the
% version of 2026-11-19 (commit b27e52b). See README.md in this folder.
\documentclass[aps,pra,onecolumn,notitlepage,11pt,amsmath,amssymb,nofootinbib]{revtex4-2}

\usepackage{bm}
\usepackage{array}
\usepackage[colorlinks=true,allcolors=blue]{hyperref}
\newcolumntype{R}[1]{>{\raggedright\arraybackslash}p{#1}}

\newcommand{\Nop}{\hat N}
\newcommand{\Iop}{\hat I}
\newcommand{\Qop}{\hat Q}
\newcommand{\Xop}{\hat X}
\newcommand{\dd}{\mathrm{d}}
\newcommand{\ee}{\mathrm{e}}
\newcommand{\ii}{\mathrm{i}}
\newcommand{\sinc}{\operatorname{sinc}}
\newcommand{\TN}[1]{\text{TN\,#1}}

\begin{document}

\title{Detailed derivation of Eqs.~(7)--(14) of the theory note:\\
evolution operator, input-output relations, and phase modulation of the pump}
\author{Working notes, 20 November 2026}
\noaffiliation

\begin{abstract}
Equations (7)--(14) of the theory note ``Quantum theory of optical
rectification with a quantized pump'' state the exact solution of the model:
the evolution operator, the output THz field, the conservation of the photon
flux, and the phase modulation of the pump. These notes derive each of them
from the generator, Eq.~(3) of the theory note, with every intermediate step
written out. The eight target equations carry the tags (TN\,7) to (TN\,14);
all other equation numbers are local. Equation numbers of the theory note
refer to its version of 19 November 2026. Section~\ref{sec:checks} adds three
consistency checks that are not in the theory note.
\end{abstract}

\maketitle

\section{What is derived}

\begin{center}
\begin{tabular}{lll}
\hline\hline
Tag & Statement & Derived in\\
\hline
(TN\,7) & $\hat U=\ee^{\ii\hat\Phi}\hat D$ & Sec.~\ref{sec:U}\\
(TN\,8) & $\hat D=\exp\{-\ii\int\dd\Omega\,[\zeta\,\hat b^\dagger\Nop+\mathrm{H.c.}]\}$ & Sec.~\ref{sec:first}\\
(TN\,9) & $\zeta=\kappa L\,\ee^{-\ii\Delta kL/2}\sinc(\Delta kL/2)$ & Sec.~\ref{sec:first}\\
(TN\,10) & $\hat b_{\rm out}=\hat b_{\rm in}-\ii\zeta\Nop$ & Sec.~\ref{sec:b}\\
(TN\,11) & $\Iop_{\rm out}(\tau)=\Iop_{\rm in}(\tau)$ & Sec.~\ref{sec:flux}\\
(TN\,12) & $\hat a_{\rm out}(\tau)=\ee^{-\ii\hat\vartheta(\tau)}\hat a_{\rm in}(\tau)$ & Sec.~\ref{sec:a}\\
(TN\,13) & $\hat\vartheta=\tfrac12[\hat F_{\rm in}+\hat F_{\rm out}]$ at $\Delta k=0$ & Sec.~\ref{sec:a}\\
(TN\,14) & $\hat F(\tau)=\int\dd\Omega\,[\zeta\,\hat b^\dagger\ee^{\ii\Omega\tau}+\mathrm{H.c.}]$ & Secs.~\ref{sec:time}, \ref{sec:a}\\
\hline\hline
\end{tabular}
\end{center}

The phase operator $\hat\Phi$ that appears in (TN\,7) is given in Appendix~A
of the theory note; it is derived here in Sec.~\ref{sec:second}.

\section{Starting point}
\label{sec:start}

\emph{Operators} [Eqs.~(1) and (2) of the theory note]. The pump is described
by an envelope operator in the frame moving with its group velocity, and the
THz field by continuum mode operators:
\begin{equation}
[\hat a(\tau),\hat a^\dagger(\tau')]=\delta(\tau-\tau'),\qquad
[\hat b(\Omega),\hat b^\dagger(\Omega')]=\delta(\Omega-\Omega')\quad(\Omega>0),
\end{equation}
and every pump operator commutes with every THz operator. The photon flux and
its Fourier transform are
\begin{equation}
\Iop(\tau)=\hat a^\dagger(\tau)\hat a(\tau),\qquad
\Nop(\Omega)=\int\!\dd\tau\,\ee^{\ii\Omega\tau}\,\Iop(\tau),\qquad
\Nop^\dagger(\Omega)=\int\!\dd\tau\,\ee^{-\ii\Omega\tau}\,\Iop(\tau)=\Nop(-\Omega).
\label{eq:Ndef}
\end{equation}

\emph{Generator} [Eq.~(3) of the theory note]. With the abbreviation
$\Xop(\Omega)=\hat b^\dagger(\Omega)\Nop(\Omega)$,
\begin{equation}
\hat G(z)=\hbar\int_0^\infty\!\!\dd\Omega\;\kappa(\Omega)
\Big[\ee^{-\ii\Delta k\,z}\,\Xop(\Omega)+\ee^{\ii\Delta k\,z}\,\Xop^\dagger(\Omega)\Big],
\qquad \Delta k(\Omega)=k_T(\Omega)-\frac{\Omega}{v_g},
\label{eq:G}
\end{equation}
with $\kappa(\Omega)$ real. $\Xop$ creates a THz photon at $\Omega$ and shifts a
pump photon down by $\Omega$; $\Xop^\dagger=\hat b\,\Nop^\dagger$ does the
reverse.

\emph{Evolution.} The evolution operator through a crystal of length $L$ is
the solution of
\begin{equation}
\frac{\partial\hat U(z)}{\partial z}=-\frac{\ii}{\hbar}\,\hat G(z)\,\hat U(z),
\qquad \hat U(0)=1,\qquad \hat U\equiv\hat U(L),
\label{eq:evol}
\end{equation}
and for any operator $\hat O_{\rm out}=\hat U^\dagger\hat O_{\rm in}\hat U$.
``In'' operators are the operators above; the subscript is dropped where no
confusion can arise.

\emph{Assumption on $\zeta$.} The integrals $\int\dd\Omega\,|\zeta|^2$ and
$\int\dd\Omega\,\Omega|\zeta|^2$ are taken to be finite, which holds for any
finite phase-matching or detection bandwidth.

\section{Commutators}
\label{sec:comm}

\emph{(L1) The flux commutes with itself} [Eq.~(5) of the theory note].
Moving $\hat a(\tau)$ to the right of $\hat a^\dagger(\tau')$,
\begin{equation}
\Iop(\tau)\Iop(\tau')=\hat a^\dagger(\tau)\big[\hat a^\dagger(\tau')\hat a(\tau)
+\delta(\tau-\tau')\big]\hat a(\tau')
=\hat a^\dagger(\tau)\hat a^\dagger(\tau')\hat a(\tau)\hat a(\tau')
+\delta(\tau-\tau')\,\hat a^\dagger(\tau)\hat a(\tau').
\label{eq:II}
\end{equation}
The first term is symmetric under $\tau\leftrightarrow\tau'$. In the second,
the delta function allows $\tau'$ to be replaced by $\tau$, so it equals
$\delta(\tau-\tau')\Iop(\tau)$ and is symmetric too. Hence
\begin{equation}
[\Iop(\tau),\Iop(\tau')]=0,\qquad
\Iop(\tau)\Iop(\tau')={:}\Iop(\tau)\Iop(\tau'){:}+\delta(\tau-\tau')\Iop(\tau).
\label{eq:L1}
\end{equation}

\emph{(L2) Flux and envelope.} From
$\hat a(\tau')\hat a(\tau)=\hat a(\tau)\hat a(\tau')$ and
$\hat a^\dagger(\tau')\hat a(\tau)=\hat a(\tau)\hat a^\dagger(\tau')-\delta(\tau-\tau')$,
\begin{equation}
\Iop(\tau')\hat a(\tau)=\hat a^\dagger(\tau')\hat a(\tau)\hat a(\tau')
=\hat a(\tau)\Iop(\tau')-\delta(\tau-\tau')\hat a(\tau')
\;\;\Longrightarrow\;\;
[\Iop(\tau'),\hat a(\tau)]=-\delta(\tau-\tau')\,\hat a(\tau).
\label{eq:L2}
\end{equation}
The adjoint is
$[\Iop(\tau'),\hat a^\dagger(\tau)]=+\delta(\tau-\tau')\hat a^\dagger(\tau)$.
For a linear functional of the flux this gives
\begin{equation}
\hat\Lambda=\int\!\dd\tau'\,\lambda(\tau')\Iop(\tau')+\text{const}:\qquad
[\hat\Lambda,\hat a(\tau)]=-\lambda(\tau)\hat a(\tau),\qquad
\hat a(\tau)\hat\Lambda=[\hat\Lambda+\lambda(\tau)]\,\hat a(\tau),
\end{equation}
and therefore for any function $f$
\begin{equation}
\hat a(\tau)\,f(\hat\Lambda)=f\big(\hat\Lambda+\lambda(\tau)\big)\,\hat a(\tau),
\qquad
f(\hat\Lambda)\,\hat a^\dagger(\tau)=\hat a^\dagger(\tau)\,f\big(\hat\Lambda+\lambda(\tau)\big).
\label{eq:shift}
\end{equation}
Annihilating a photon at $\tau$ lowers the flux by $\delta(\tau'-\tau)$.

\emph{(L3) Fourier components.} Multiplying Eq.~\eqref{eq:L2} by
$\ee^{\pm\ii\Omega\tau'}$ and integrating over $\tau'$,
\begin{equation}
[\Nop(\Omega),\hat a(\tau)]=-\ee^{\ii\Omega\tau}\,\hat a(\tau),\qquad
[\Nop^\dagger(\Omega),\hat a(\tau)]=-\ee^{-\ii\Omega\tau}\,\hat a(\tau).
\label{eq:L3}
\end{equation}

\emph{(L4)} From (L1) and Eq.~\eqref{eq:Ndef},
\begin{equation}
[\Nop(\Omega),\Nop(\Omega')]=[\Nop(\Omega),\Nop^\dagger(\Omega')]=0 .
\label{eq:L4}
\end{equation}

\emph{(L5) The two elementary processes.} Because all $\hat b^\dagger$ commute
among themselves and so do all $\Nop$,
$[\Xop(\Omega),\Xop(\Omega')]=0$. For the mixed commutator, using (L4) to move
the flux operators freely,
\begin{align}
[\Xop(\Omega),\Xop^\dagger(\Omega')]
&=\hat b^\dagger(\Omega)\Nop(\Omega)\,\hat b(\Omega')\Nop^\dagger(\Omega')
-\hat b(\Omega')\Nop^\dagger(\Omega')\,\hat b^\dagger(\Omega)\Nop(\Omega)\nonumber\\
&=\big[\hat b^\dagger(\Omega),\hat b(\Omega')\big]\,\Nop(\Omega)\Nop^\dagger(\Omega')
=-\delta(\Omega-\Omega')\,\Qop(\Omega),\qquad
\Qop(\Omega)\equiv\Nop^\dagger(\Omega)\Nop(\Omega).
\label{eq:L5}
\end{align}
In a model where $\Nop$ and $\Nop^\dagger$ did not commute, an extra term
$\hat b(\Omega')\hat b^\dagger(\Omega)[\Nop(\Omega),\Nop^\dagger(\Omega')]$
would appear here, and the solution below would not close.

\emph{(L6) Generator at two positions} [Eq.~(6) of the theory note].
Inserting Eq.~\eqref{eq:G} twice and using (L5),
\begin{align}
[\hat G(z),\hat G(z')]
&=\hbar^2\!\iint\!\dd\Omega\,\dd\Omega'\,\kappa\kappa'
\Big\{\ee^{-\ii\Delta kz+\ii\Delta k'z'}[\Xop(\Omega),\Xop^\dagger(\Omega')]\nonumber\\
&\hspace{9em}+\ee^{\ii\Delta kz-\ii\Delta k'z'}[\Xop^\dagger(\Omega),\Xop(\Omega')]\Big\}\nonumber\\
&=\hbar^2\!\int\!\dd\Omega\,\kappa^2
\Big\{-\ee^{-\ii\Delta k(z-z')}+\ee^{\ii\Delta k(z-z')}\Big\}\Qop(\Omega)
=2\ii\hbar^2\!\int_0^\infty\!\!\dd\Omega\,\kappa^2
\sin[\Delta k\,(z-z')]\;\Qop(\Omega).
\label{eq:L6}
\end{align}
Primes on $\kappa$ and $\Delta k$ denote the argument $\Omega'$. The terms
with $[\Xop,\Xop']$ and $[\Xop^\dagger,\Xop'^\dagger]$ vanish. The second
commutator shown is $[\Xop^\dagger(\Omega),\Xop(\Omega')]
=+\delta(\Omega-\Omega')\Qop(\Omega)$.

\emph{(L7)} $\Qop(\Omega)$ is built from flux operators only, so by (L4) it
commutes with $\Xop$ and $\Xop^\dagger$, hence with $\hat G(z)$ for all $z$:
\begin{equation}
[\Qop(\Omega),\hat G(z)]=0,\qquad
\big[[\hat G(z),\hat G(z')],\hat G(z'')\big]=0 .
\label{eq:L7}
\end{equation}

\emph{(L8) Field commutator.} For the Hermitian THz operator of (TN\,14),
$\hat F(\tau)=\int\dd\Omega\,[\zeta\,\hat b^\dagger\ee^{\ii\Omega\tau}
+\zeta^*\hat b\,\ee^{-\ii\Omega\tau}]$,
\begin{align}
[\hat F(\tau'),\hat F(\tau)]
&=\iint\!\dd\Omega\,\dd\Omega'\Big\{\zeta\zeta'^*\,
\ee^{\ii\Omega\tau'-\ii\Omega'\tau}\,[\hat b^\dagger(\Omega),\hat b(\Omega')]
+\zeta^*\zeta'\,\ee^{-\ii\Omega\tau'+\ii\Omega'\tau}\,[\hat b(\Omega),\hat b^\dagger(\Omega')]\Big\}\nonumber\\
&=\int\!\dd\Omega\,|\zeta|^2\Big\{-\ee^{\ii\Omega(\tau'-\tau)}+\ee^{-\ii\Omega(\tau'-\tau)}\Big\}
=-2\ii\,S(\tau'-\tau),
\label{eq:L8}
\end{align}
\begin{equation}
S(x)\equiv\int_0^\infty\!\!\dd\Omega\,|\zeta(\Omega)|^2\sin\Omega x .
\label{eq:S}
\end{equation}
$S$ is real and odd, $S(0)=0$, and $S'(0)=\int\dd\Omega\,\Omega|\zeta|^2$.

\section{Operator identities}
\label{sec:ident}

\emph{(I1) Conjugation.} For operators $\hat A$, $\hat B$,
\begin{equation}
\ee^{\hat A}\hat B\,\ee^{-\hat A}=\hat B+[\hat A,\hat B]
+\frac{1}{2!}[\hat A,[\hat A,\hat B]]+\dots
\label{eq:I1}
\end{equation}
If $[\hat A,[\hat A,\hat B]]=0$ the series ends after the second term.

\emph{(I2) Derivative of an exponential.} For an operator $\hat M(s)$ that
depends on a parameter, with $\hat M'=\partial_s\hat M$,
\begin{equation}
\frac{\partial}{\partial s}\ee^{\hat M}
=\ee^{\hat M}\!\int_0^1\!\dd\lambda\;\ee^{-\lambda\hat M}\,\hat M'\,\ee^{\lambda\hat M}.
\label{eq:duhamel}
\end{equation}
Proof: $\hat K(\lambda)=\ee^{-\lambda\hat M}\partial_s\ee^{\lambda\hat M}$
obeys $\hat K(0)=0$ and
$\partial_\lambda\hat K=-\ee^{-\lambda\hat M}\hat M\partial_s\ee^{\lambda\hat M}
+\ee^{-\lambda\hat M}\partial_s(\hat M\ee^{\lambda\hat M})
=\ee^{-\lambda\hat M}\hat M'\ee^{\lambda\hat M}$; integrate from $0$ to $1$.
If $[\hat M,[\hat M,\hat M']]=0$, (I1) gives
$\ee^{-\lambda\hat M}\hat M'\ee^{\lambda\hat M}=\hat M'-\lambda[\hat M,\hat M']$,
the integral is $\hat M'-\tfrac12[\hat M,\hat M']$, and moving $\ee^{\hat M}$
to the right with $\ee^{\hat M}\hat M'=(\hat M'+[\hat M,\hat M'])\ee^{\hat M}$,
\begin{equation}
\frac{\partial}{\partial s}\ee^{\hat M}
=\Big(\hat M'+\tfrac12[\hat M,\hat M']\Big)\,\ee^{\hat M}
\qquad\text{if }[\hat M,[\hat M,\hat M']]=0 .
\label{eq:I2}
\end{equation}

\emph{(I3) Reordering.} If $[\hat A,\hat B]$ commutes with $\hat A$ and
$\hat B$, then $\ee^{\hat A}\ee^{\hat B}=\ee^{\hat B}\ee^{\hat A}\,\ee^{[\hat A,\hat B]}$.

\section{Evolution operator: Eqs.~(7)--(9)}
\label{sec:U}

\subsection{Form of the solution}

The Magnus expansion writes the solution of Eq.~\eqref{eq:evol} as a single
exponential $\exp(\hat\Omega_1+\hat\Omega_2+\hat\Omega_3+\dots)$, where
$\hat\Omega_1$ is the integral of the generator, $\hat\Omega_2$ contains one
commutator of the generator at two positions, and every $\hat\Omega_n$ with
$n\ge3$ is a sum of nested commutators that contain $[[\hat G,\hat G],\hat G]$
\cite{blanes2009}. By (L7) these vanish. This suggests the two-factor form
\begin{equation}
\hat U(z)=\ee^{\hat\Omega_2(z)}\,\ee^{\hat\Omega_1(z)},
\label{eq:ansatz}
\end{equation}
\begin{equation}
\hat\Omega_1(z)=-\frac{\ii}{\hbar}\int_0^z\!\dd z_1\,\hat G(z_1),\qquad
\hat\Omega_2(z)=-\frac{1}{2\hbar^2}\int_0^z\!\dd z_1\!\int_0^{z_1}\!\!\dd z_2\,
[\hat G(z_1),\hat G(z_2)] .
\label{eq:Om}
\end{equation}
Instead of relying on the expansion we verify Eq.~\eqref{eq:ansatz} directly.

\subsection{Proof by differentiation}

By (L6) and (L7), $\hat\Omega_2(z)$ is a function of the operators $\Qop$ only
and commutes with $\hat G$ at any position, hence with $\hat\Omega_1(z)$ and
with $\hat\Omega_1'(z)=-(\ii/\hbar)\hat G(z)$. The commutator
\begin{equation}
[\hat\Omega_1(z),\hat\Omega_1'(z)]
=\Big(-\frac{\ii}{\hbar}\Big)^{2}\!\int_0^z\!\dd z_2\,[\hat G(z_2),\hat G(z)]
=\frac{1}{\hbar^2}\int_0^z\!\dd z_2\,[\hat G(z),\hat G(z_2)]
\label{eq:O1O1p}
\end{equation}
is also a function of the $\Qop$ only, so the condition of (I2) holds and
\begin{equation}
\frac{\partial}{\partial z}\ee^{\hat\Omega_1}
=\Big(\hat\Omega_1'+\tfrac12[\hat\Omega_1,\hat\Omega_1']\Big)\ee^{\hat\Omega_1}.
\end{equation}
Since $\hat\Omega_2$ commutes with its own derivative,
$\partial_z\ee^{\hat\Omega_2}=\hat\Omega_2'\ee^{\hat\Omega_2}$, and
\begin{equation}
\frac{\partial\hat U}{\partial z}
=\Big(\hat\Omega_2'+\hat\Omega_1'+\tfrac12[\hat\Omega_1,\hat\Omega_1']\Big)\hat U .
\end{equation}
From Eq.~\eqref{eq:Om},
$\hat\Omega_2'(z)=-(2\hbar^2)^{-1}\int_0^z\dd z_2\,[\hat G(z),\hat G(z_2)]$,
which by Eq.~\eqref{eq:O1O1p} is exactly
$-\tfrac12[\hat\Omega_1,\hat\Omega_1']$. The two terms cancel and
\begin{equation}
\frac{\partial\hat U}{\partial z}=\hat\Omega_1'(z)\,\hat U
=-\frac{\ii}{\hbar}\hat G(z)\,\hat U,\qquad \hat U(0)=1 .
\end{equation}
Equation~\eqref{eq:ansatz} therefore solves Eq.~\eqref{eq:evol}, and the
solution of a first-order linear equation with a given initial value is
unique. \emph{No expansion in $\kappa$ has been made.}

The role of the two factors: $\ee^{\hat\Omega_1}$ alone would be the solution
if the generators at different positions commuted. They do not, because THz
radiation emitted at one position is present at the next; the correction for
this is $\ee^{\hat\Omega_2}$. There is no third factor because the correction
involves the flux only, and the flux is not changed by the interaction.

\subsection{First factor: \texorpdfstring{$\hat D$}{D} and \texorpdfstring{$\zeta$}{zeta}}
\label{sec:first}

Inserting Eq.~\eqref{eq:G} in $\hat\Omega_1(L)$,
\begin{equation}
\hat\Omega_1(L)=-\ii\int_0^\infty\!\!\dd\Omega\;\kappa
\Big[\Xop(\Omega)\int_0^L\!\!\dd z\,\ee^{-\ii\Delta kz}
+\Xop^\dagger(\Omega)\int_0^L\!\!\dd z\,\ee^{\ii\Delta kz}\Big].
\end{equation}
The integral over the crystal is
\begin{equation}
\int_0^L\!\!\dd z\,\ee^{-\ii\Delta kz}=\frac{1-\ee^{-\ii\Delta kL}}{\ii\Delta k}
=\ee^{-\ii\Delta kL/2}\,\frac{\ee^{\ii\Delta kL/2}-\ee^{-\ii\Delta kL/2}}{\ii\Delta k}
=L\,\ee^{-\ii\Delta kL/2}\,\frac{\sin(\Delta kL/2)}{\Delta kL/2},
\end{equation}
and the second integral is its complex conjugate. With
\begin{equation}
\boxed{\;\zeta(\Omega)=\kappa(\Omega)\,L\,\ee^{-\ii\Delta kL/2}\,
\sinc\!\Big(\frac{\Delta kL}{2}\Big),\qquad\sinc x=\frac{\sin x}{x}\;}
\tag{\TN{9}}
\end{equation}
the first factor is
\begin{equation}
\boxed{\;\hat D\equiv\ee^{\hat\Omega_1(L)}
=\exp\!\Big\{-\ii\int_0^\infty\!\!\dd\Omega\,
\big[\zeta(\Omega)\,\hat b^\dagger(\Omega)\Nop(\Omega)
+\zeta^*(\Omega)\,\hat b(\Omega)\Nop^\dagger(\Omega)\big]\Big\}\;}
\tag{\TN{8}}
\end{equation}
The exponent is $-\ii$ times a Hermitian operator, so $\hat D$ is unitary.
For perfect phase matching $\zeta=\kappa L$; for a large mismatch
$|\zeta|\le2\kappa/|\Delta k|$ and stops growing with $L$.

\subsection{Second factor: \texorpdfstring{$\hat\Phi$}{Phi}}
\label{sec:second}

Inserting (L6) in $\hat\Omega_2(L)$,
\begin{equation}
\hat\Omega_2(L)=-\frac{1}{2\hbar^2}\cdot2\ii\hbar^2\!\int\!\dd\Omega\,\kappa^2\,\Qop(\Omega)
\int_0^L\!\!\dd z_1\!\int_0^{z_1}\!\!\dd z_2\,\sin[\Delta k(z_1-z_2)] .
\end{equation}
The inner integral is $[1-\cos(\Delta kz_1)]/\Delta k$, and the outer one
\begin{equation}
\int_0^L\!\!\dd z_1\,\frac{1-\cos\Delta kz_1}{\Delta k}
=\frac{1}{\Delta k}\Big[L-\frac{\sin\Delta kL}{\Delta k}\Big]
=\frac{\Delta kL-\sin\Delta kL}{\Delta k^2}.
\end{equation}
Hence $\hat\Omega_2(L)=\ii\hat\Phi$ with
\begin{equation}
\hat\Phi=-\int_0^\infty\!\!\dd\Omega\;\varphi(\Omega)\,\Nop^\dagger(\Omega)\Nop(\Omega),
\qquad
\varphi(\Omega)=\kappa^2\,\frac{\Delta kL-\sin\Delta kL}{\Delta k^2},
\label{eq:Phi}
\end{equation}
which is the expression in Appendix~A of the theory note. $\varphi$ is real
and $\Qop$ is Hermitian, so $\hat\Phi$ is Hermitian and $\ee^{\ii\hat\Phi}$ is
unitary. For $\Delta kL\ll1$, $\sin x=x-x^3/6+\dots$ gives
$\varphi\simeq\kappa^2\Delta kL^3/6\to0$; for $\Delta kL\gg1$,
$\varphi\simeq\kappa^2L/\Delta k$. Together with Eq.~\eqref{eq:ansatz},
\begin{equation}
\boxed{\;\hat U=\ee^{\ii\hat\Phi}\,\hat D\;}
\tag{\TN{7}}
\end{equation}
The two factors commute, so the order is immaterial.

\subsection{Time-domain forms}
\label{sec:time}

Both exponents can be written with the flux $\Iop(\tau)$. Inserting
Eq.~\eqref{eq:Ndef} in the exponent of (TN\,8) and exchanging the integrals,
\begin{equation}
\int\!\dd\Omega\,\big[\zeta\,\hat b^\dagger\Nop+\zeta^*\hat b\,\Nop^\dagger\big]
=\int\!\dd\tau\;\Iop(\tau)\int\!\dd\Omega\,
\big[\zeta\,\hat b^\dagger\ee^{\ii\Omega\tau}+\zeta^*\hat b\,\ee^{-\ii\Omega\tau}\big]
=\int\!\dd\tau\;\Iop(\tau)\,\hat F(\tau),
\label{eq:Y}
\end{equation}
with the Hermitian THz operator
\begin{equation}
\boxed{\;\hat F(\tau)=\int_0^\infty\!\!\dd\Omega\,
\big[\zeta(\Omega)\,\hat b^\dagger(\Omega)\,\ee^{\ii\Omega\tau}+\mathrm{H.c.}\big]\;}
\tag{\TN{14}}
\end{equation}
so that $\hat D=\exp[-\ii\int\dd\tau\,\Iop(\tau)\hat F(\tau)]$: every photon
that passes at time $\tau$ contributes the phase operator $\hat F(\tau)$.
$\hat F$ is dimensionless; it is the THz field in the units in which it acts
as a phase on the pump. Its vacuum variance is
$\langle\hat F^2\rangle=\int\dd\Omega\,|\zeta|^2\equiv B_T$.

For $\hat\Phi$, Eq.~\eqref{eq:Ndef} gives
$\Qop(\Omega)=\iint\dd\tau\,\dd\tau'\,\ee^{-\ii\Omega(\tau-\tau')}\Iop(\tau)\Iop(\tau')$.
The product $\Iop(\tau)\Iop(\tau')$ is symmetric by (L1), so only the even
part of the exponential survives:
\begin{equation}
\hat\Phi=-\iint\!\dd\tau\,\dd\tau'\;V(\tau-\tau')\,\Iop(\tau)\Iop(\tau'),\qquad
V(x)=\int_0^\infty\!\!\dd\Omega\,\varphi(\Omega)\cos\Omega x .
\label{eq:PhiV}
\end{equation}
This is an interaction between the flux at two times, mediated by the THz
field.

\section{Output THz field: Eq.~(10)}
\label{sec:b}

\emph{From the evolution operator.} $\hat\Phi$ contains pump operators only
and commutes with $\hat b$, so
$\hat b_{\rm out}=\hat D^\dagger\ee^{-\ii\hat\Phi}\hat b\,\ee^{\ii\hat\Phi}\hat D
=\hat D^\dagger\hat b\hat D$. Write $\hat D=\ee^{-\ii\hat Y}$ with
$\hat Y=\int\dd\Omega'\,[\zeta'\hat b'^\dagger\Nop'+\zeta'^*\hat b'\Nop'^\dagger]$,
and apply (I1) with $\hat A=\ii\hat Y$:
\begin{equation}
[\hat Y,\hat b(\Omega)]=\int\!\dd\Omega'\,\zeta'\,\Nop'\,[\hat b^\dagger(\Omega'),\hat b(\Omega)]
=-\zeta(\Omega)\Nop(\Omega),\qquad
[\hat Y,[\hat Y,\hat b(\Omega)]]=-\zeta(\Omega)[\hat Y,\Nop(\Omega)]=0 .
\end{equation}
The second commutator vanishes by (L4). The series ends after the first
commutator:
\begin{equation}
\boxed{\;\hat b_{\rm out}(\Omega)=\hat b(\Omega)+\ii[\hat Y,\hat b(\Omega)]
=\hat b_{\rm in}(\Omega)-\ii\,\zeta(\Omega)\,\Nop(\Omega)\;}
\tag{\TN{10}}
\end{equation}
This is exact; it is the first-order result of perturbation theory because
all higher orders are identically zero.

\emph{From the equation of motion.} With
$\hat b(\Omega,z)=\hat U^\dagger(z)\hat b(\Omega)\hat U(z)$ and
Eq.~\eqref{eq:evol},
\begin{equation}
\frac{\partial\hat b(\Omega,z)}{\partial z}
=\frac{\ii}{\hbar}\,\hat U^\dagger(z)\,[\hat G(z),\hat b(\Omega)]\,\hat U(z),
\qquad
[\hat G(z),\hat b(\Omega)]=-\hbar\kappa(\Omega)\,\ee^{-\ii\Delta kz}\,\Nop(\Omega).
\end{equation}
Because $[\hat G(z),\Nop]=0$, $\hat U^\dagger(z)\Nop\hat U(z)=\Nop$ for all
$z$, and
\begin{equation}
\frac{\partial\hat b(\Omega,z)}{\partial z}=-\ii\kappa\,\ee^{-\ii\Delta kz}\,\Nop(\Omega)
\quad\Longrightarrow\quad
\hat b(\Omega,L)=\hat b(\Omega)-\ii\kappa\int_0^L\!\!\dd z\,\ee^{-\ii\Delta kz}\,\Nop(\Omega)
=\hat b(\Omega)-\ii\zeta\Nop(\Omega).
\end{equation}
The source of the THz field does not evolve, so the equation is integrated as
it stands. The THz amplitude grows linearly along the crystal for
$\Delta k=0$.

\section{Conservation of the flux: Eq.~(11)}
\label{sec:flux}

$\hat\Phi$ is a function of the flux, Eq.~\eqref{eq:PhiV}, and the exponent of
$\hat D$ is the flux times a THz operator, Eq.~\eqref{eq:Y}. By (L1) both
commute with $\Iop(\tau)$ for every $\tau$, so
\begin{equation}
\boxed{\;\Iop_{\rm out}(\tau)=\hat U^\dagger\Iop(\tau)\hat U=\Iop_{\rm in}(\tau)\;}
\tag{\TN{11}}
\end{equation}
The same holds for $\Nop(\Omega)$ at every $\Omega$, for the total photon
number $\Nop(0)$, and for every function of the flux: all intensity
correlation functions of the pump are unchanged by the crystal.

\section{Output pump field: Eqs.~(12)--(14)}
\label{sec:a}

The output is
$\hat a_{\rm out}(\tau)=\hat D^\dagger\big[\ee^{-\ii\hat\Phi}\hat a(\tau)\ee^{\ii\hat\Phi}\big]\hat D$.
The two conjugations are done in turn.

\subsection{Conjugation with \texorpdfstring{$\ee^{\ii\hat\Phi}$}{exp(i Phi)}}

From Eq.~\eqref{eq:PhiV} and the product rule,
\begin{equation}
[\Iop(\tau_1)\Iop(\tau_2),\hat a(\tau)]
=\Iop(\tau_1)[\Iop(\tau_2),\hat a(\tau)]+[\Iop(\tau_1),\hat a(\tau)]\Iop(\tau_2)
=-\delta(\tau_2-\tau)\Iop(\tau_1)\hat a(\tau)-\delta(\tau_1-\tau)\hat a(\tau)\Iop(\tau_2).
\end{equation}
In the last term $\hat a(\tau)$ is moved to the right with (L2),
$\hat a(\tau)\Iop(\tau_2)=\Iop(\tau_2)\hat a(\tau)+\delta(\tau-\tau_2)\hat a(\tau)$.
Integrating with $-V(\tau_1-\tau_2)$ and using that $V$ is even,
\begin{equation}
[\hat\Phi,\hat a(\tau)]=\hat\theta(\tau)\,\hat a(\tau),\qquad
\hat\theta(\tau)=2\int\!\dd\tau'\,V(\tau-\tau')\,\Iop(\tau')+V(0).
\label{eq:thetaK}
\end{equation}
$\hat\theta(\tau)$ is a function of the flux and commutes with $\hat\Phi$, so
$[\hat\Phi,[\hat\Phi,\hat a]]=\hat\theta^2\hat a$ and so on, and (I1) with
$\hat A=-\ii\hat\Phi$ sums to
\begin{equation}
\ee^{-\ii\hat\Phi}\,\hat a(\tau)\,\ee^{\ii\hat\Phi}
=\sum_{n=0}^\infty\frac{(-\ii)^n}{n!}\,\hat\theta(\tau)^n\,\hat a(\tau)
=\ee^{-\ii\hat\theta(\tau)}\,\hat a(\tau).
\label{eq:step1}
\end{equation}
The constant $V(0)$ is the phase a single photon acquires by emitting and
reabsorbing THz radiation. Check with one photon at $\tau_0$:
$\hat\Phi|\tau_0\rangle=-V(0)|\tau_0\rangle$ and $\hat\Phi|0\rangle=0$, so the
left side applied to $|\tau_0\rangle$ gives
$\ee^{-\ii V(0)}\delta(\tau-\tau_0)|0\rangle$; on the right,
$\hat\theta(\tau)|0\rangle=V(0)|0\rangle$ gives the same.

\subsection{Conjugation with \texorpdfstring{$\hat D$}{D}}

$\hat\theta(\tau)$ commutes with $\hat D$ (Sec.~\ref{sec:flux}), so
$\hat D^\dagger\ee^{-\ii\hat\theta}\hat a\hat D
=\ee^{-\ii\hat\theta}\hat D^\dagger\hat a\hat D$. With
$\hat D=\ee^{-\ii\hat Y}$ and $\hat Y=\int\dd\tau'\,\Iop(\tau')\hat F(\tau')$,
(L2) gives
\begin{equation}
[\hat Y,\hat a(\tau)]=\int\!\dd\tau'\,\hat F(\tau')\,[\Iop(\tau'),\hat a(\tau)]
=-\hat F(\tau)\,\hat a(\tau).
\label{eq:Ya}
\end{equation}
Unlike in Sec.~\ref{sec:b}, the series (I1) does not end here, because
$\hat F(\tau)$ does not commute with $\hat Y$: by (L8),
\begin{equation}
[\hat Y,\hat F(\tau)]=\int\!\dd\tau'\,\Iop(\tau')\,[\hat F(\tau'),\hat F(\tau)]
=-2\ii\,\hat\Sigma(\tau),\qquad
\hat\Sigma(\tau)\equiv\int\!\dd\tau'\,S(\tau'-\tau)\,\Iop(\tau').
\label{eq:YF}
\end{equation}
The series is summed with a differential equation. Let
$\hat A(\lambda)=\ee^{\ii\lambda\hat Y}\hat a(\tau)\ee^{-\ii\lambda\hat Y}$, so
that $\hat A(0)=\hat a(\tau)$ and $\hat A(1)=\hat D^\dagger\hat a(\tau)\hat D$.
With Eq.~\eqref{eq:Ya},
\begin{equation}
\frac{\partial\hat A}{\partial\lambda}
=\ii\,\ee^{\ii\lambda\hat Y}[\hat Y,\hat a(\tau)]\ee^{-\ii\lambda\hat Y}
=-\ii\,\hat F_\lambda(\tau)\,\hat A(\lambda),\qquad
\hat F_\lambda(\tau)\equiv\ee^{\ii\lambda\hat Y}\hat F(\tau)\ee^{-\ii\lambda\hat Y}.
\end{equation}
For $\hat F_\lambda$ the series (I1) does end: $\hat\Sigma$ is a function of
the flux and commutes with $\hat Y$, so with Eq.~\eqref{eq:YF}
\begin{equation}
\hat F_\lambda(\tau)=\hat F(\tau)+\ii\lambda[\hat Y,\hat F(\tau)]
=\hat F(\tau)+2\lambda\,\hat\Sigma(\tau).
\label{eq:Flam}
\end{equation}
This is the THz field seen by the pump after a fraction $\lambda$ of the
interaction: the input field plus the part generated so far. The operators
$\hat F_\lambda(\tau)$ at different $\lambda$ commute with each other
($\hat F(\tau)$ commutes with itself and with $\hat\Sigma$), so the equation
for $\hat A$ is solved by an ordinary exponential,
\begin{equation}
\hat A(\lambda)=\exp\!\Big\{-\ii\int_0^\lambda\!\!\dd s\,\hat F_s(\tau)\Big\}\hat a(\tau)
=\exp\!\big\{-\ii\big[\lambda\hat F(\tau)+\lambda^2\hat\Sigma(\tau)\big]\big\}\,\hat a(\tau),
\end{equation}
as is checked by differentiation. At $\lambda=1$,
\begin{equation}
\hat D^\dagger\,\hat a(\tau)\,\hat D
=\exp\!\big\{-\ii\big[\hat F_{\rm in}(\tau)+\hat\Sigma(\tau)\big]\big\}\,\hat a(\tau).
\label{eq:step2}
\end{equation}
$\hat\Sigma(\tau)$ stands to the left of $\hat a(\tau)$ and contains the flux
at $\tau'=\tau$; moving it to the right would add $S(0)=0$ by
Eq.~\eqref{eq:shift}, so there is no ordering constant.

\subsection{The flux-dependent phase is half the generated field}

Define $\hat F_{\rm out}(\tau)=\hat U^\dagger\hat F_{\rm in}(\tau)\hat U$, which
is (TN\,14) with $\hat b_{\rm out}$ in place of $\hat b_{\rm in}$. Inserting
(TN\,10) and its adjoint
$\hat b_{\rm out}^\dagger=\hat b^\dagger+\ii\zeta^*\Nop^\dagger$,
\begin{align}
\hat F_{\rm out}(\tau)-\hat F_{\rm in}(\tau)
&=\ii\int\!\dd\Omega\,|\zeta|^2\Big[\Nop^\dagger(\Omega)\,\ee^{\ii\Omega\tau}
-\Nop(\Omega)\,\ee^{-\ii\Omega\tau}\Big]\nonumber\\
&=\ii\int\!\dd\Omega\,|\zeta|^2\int\!\dd\tau'\,
\Big[\ee^{-\ii\Omega(\tau'-\tau)}-\ee^{\ii\Omega(\tau'-\tau)}\Big]\Iop(\tau')
=2\int\!\dd\tau'\,S(\tau'-\tau)\,\Iop(\tau')=2\hat\Sigma(\tau).
\label{eq:Fdiff}
\end{align}
The same follows from Eq.~\eqref{eq:Flam} at $\lambda=1$. Therefore
\begin{equation}
\hat F_{\rm in}(\tau)+\hat\Sigma(\tau)
=\tfrac12\big[\hat F_{\rm in}(\tau)+\hat F_{\rm out}(\tau)\big].
\label{eq:half}
\end{equation}
The factor $\tfrac12$ has a simple origin. By Eq.~\eqref{eq:Flam} the field
seen by the pump grows linearly during the interaction, from
$\hat F_{\rm in}$ to $\hat F_{\rm out}$, and the accumulated phase is its
average.

\subsection{Result}

Combining Eqs.~\eqref{eq:step1}, \eqref{eq:step2}, and \eqref{eq:half}; the
three exponents commute with one another, so they add:
\begin{equation}
\boxed{\;\hat a_{\rm out}(\tau)=\ee^{-\ii\hat\vartheta(\tau)}\,\hat a_{\rm in}(\tau)\;}
\tag{\TN{12}}
\end{equation}
\begin{equation}
\hat\vartheta(\tau)=\tfrac12\big[\hat F_{\rm in}(\tau)+\hat F_{\rm out}(\tau)\big]
+\hat\theta(\tau)
=\hat F_{\rm in}(\tau)+\int\!\dd\tau'\,\big[S(\tau'-\tau)+2V(\tau-\tau')\big]\Iop(\tau')+V(0).
\label{eq:varthetafull}
\end{equation}
For perfect phase matching $\varphi=0$, so $V=0$ and $\hat\theta=0$:
\begin{equation}
\boxed{\;\hat\vartheta(\tau)=\tfrac12\big[\hat F_{\rm in}(\tau)+\hat F_{\rm out}(\tau)\big]
\qquad(\Delta k=0)\;}
\tag{\TN{13}}
\end{equation}
$\hat\vartheta$ is Hermitian: $\hat F$ is Hermitian, and $S$, $V$, $\Iop$ are
real or Hermitian. Then
$\hat a_{\rm out}^\dagger\hat a_{\rm out}
=\hat a^\dagger\ee^{\ii\hat\vartheta}\ee^{-\ii\hat\vartheta}\hat a=\hat a^\dagger\hat a$,
which is (TN\,11) again.

The three parts of the phase in Eq.~\eqref{eq:varthetafull} have distinct
origins. $\hat F_{\rm in}$ is the incoming THz field; for a vacuum input it has
zero mean and variance $B_T$, a phase noise. The $S$ term is the THz field
that the pump itself has radiated, proportional to the flux. The $V$ term is
the additional flux-dependent phase that arises when the THz wave slips with
respect to the pump.

\section{Consistency checks}
\label{sec:checks}

\subsection{Commutators of the output operators}

A unitary transformation preserves commutators, so the explicit expressions
must do so too. For the THz field, from (TN\,10) and (L4),
\begin{equation}
[\hat b_{\rm out}(\Omega),\hat b_{\rm out}^\dagger(\Omega')]
=[\hat b,\hat b'^\dagger]+\zeta\zeta'^*[\Nop(\Omega),\Nop^\dagger(\Omega')]
=\delta(\Omega-\Omega').
\end{equation}
For the pump the check is not trivial and tests the coefficient of the
flux-dependent phase. Write $\hat\vartheta(\tau)=\hat F_{\rm in}(\tau)+\hat\Lambda_\tau$,
where $\hat\Lambda_\tau$ is the linear functional of the flux in
Eq.~\eqref{eq:varthetafull}, with kernel
$\lambda_\tau(\tau')=S(\tau'-\tau)+2V(\tau-\tau')$. Then
\begin{equation}
\hat a_{\rm out}(\tau)\hat a_{\rm out}^\dagger(\tau')
=\ee^{-\ii\hat\vartheta(\tau)}\big[\delta(\tau-\tau')+\hat a^\dagger(\tau')\hat a(\tau)\big]\ee^{\ii\hat\vartheta(\tau')}.
\end{equation}
The first term is $\delta(\tau-\tau')$, since the two phases cancel at equal
times. In the second, the phases are moved to the middle with
Eq.~\eqref{eq:shift}:
\begin{equation}
\ee^{-\ii\hat\vartheta(\tau)}\hat a^\dagger(\tau')
=\hat a^\dagger(\tau')\,\ee^{-\ii[\hat\vartheta(\tau)+\lambda_\tau(\tau')]},\qquad
\hat a(\tau)\,\ee^{\ii\hat\vartheta(\tau')}
=\ee^{\ii[\hat\vartheta(\tau')+\lambda_{\tau'}(\tau)]}\,\hat a(\tau).
\end{equation}
The c-number phases combine to
$\lambda_\tau(\tau')-\lambda_{\tau'}(\tau)=2S(\tau'-\tau)$, because $S$ is odd
and $V$ is even. The two operator phases are exchanged with (I3) and (L8):
$[\hat\vartheta(\tau),\hat\vartheta(\tau')]=[\hat F(\tau),\hat F(\tau')]
=2\ii S(\tau'-\tau)$, so
$\ee^{-\ii\hat\vartheta(\tau)}\ee^{\ii\hat\vartheta(\tau')}
=\ee^{\ii\hat\vartheta(\tau')}\ee^{-\ii\hat\vartheta(\tau)}\,\ee^{2\ii S(\tau'-\tau)}$.
Altogether
\begin{equation}
\hat a_{\rm out}(\tau)\hat a_{\rm out}^\dagger(\tau')
=\delta(\tau-\tau')+\hat a^\dagger(\tau')\,\ee^{\ii\hat\vartheta(\tau')}\ee^{-\ii\hat\vartheta(\tau)}\,\hat a(\tau)\;
\ee^{2\ii S(\tau'-\tau)}\,\ee^{-2\ii S(\tau'-\tau)}
=\delta(\tau-\tau')+\hat a_{\rm out}^\dagger(\tau')\hat a_{\rm out}(\tau).
\end{equation}
The phase from the noncommuting field operators is cancelled by the phase
from the flux-dependent term. With any other coefficient in front of
$\hat\Sigma$ in Eq.~\eqref{eq:step2} the cancellation would fail: the noise
operator $\hat F_{\rm in}$ and the flux-dependent phase are tied to each
other by unitarity, as in the theories of phase modulation with a
non-instantaneous response \cite{boivin1994,shapirorazavi2007}.
$[\hat a_{\rm out}(\tau),\hat a_{\rm out}(\tau')]=0$ follows in the same way.

\subsection{Energy balance}

The pump energy relative to the carrier is
$\hat W=\hbar\int\dd\tau\,\hat a^\dagger(\tau)\,\ii\partial_\tau\hat a(\tau)$.
Inserting (TN\,12),
\begin{equation}
\hat W_{\rm out}-\hat W_{\rm in}
=\hbar\int\!\dd\tau\;\hat a^\dagger(\tau)\Big[\ee^{\ii\hat\vartheta}\,\ii\partial_\tau\ee^{-\ii\hat\vartheta}\Big]\hat a(\tau).
\end{equation}
By (I2) with $\hat M=-\ii\hat\vartheta(\tau)$ and $\tau$ as the parameter; the
commutator of $\hat\vartheta$ with its derivative is a number, so the
condition holds:
$\partial_\tau\ee^{-\ii\hat\vartheta}
=(-\ii\hat\vartheta'-\tfrac12[\hat\vartheta,\hat\vartheta'])\ee^{-\ii\hat\vartheta}$,
and with (I1),
\begin{equation}
\ee^{\ii\hat\vartheta}\,\ii\partial_\tau\ee^{-\ii\hat\vartheta}
=\hat\vartheta'+\tfrac{\ii}{2}[\hat\vartheta,\hat\vartheta'],\qquad
[\hat\vartheta(\tau),\hat\vartheta'(\tau)]
=\partial_{\tau'}[\hat F(\tau),\hat F(\tau')]\big|_{\tau'=\tau}=2\ii\,S'(0).
\end{equation}
Hence the bracket equals $\hat\vartheta'(\tau)-S'(0)$. In
$\hat\vartheta'=\hat F_{\rm in}'+\hat\Sigma'+\hat\theta'$ the three terms give:
(i) $\int\dd\tau\,\Iop\,\hat F_{\rm in}'$, the work done by the incoming THz
field, with zero mean for a vacuum input; (ii) with
$\hat\Sigma'(\tau)=-\int\dd\tau'\,S'(\tau'-\tau)\Iop(\tau')$ and
$\hat a^\dagger(\tau)\Iop(\tau')\hat a(\tau)={:}\Iop(\tau)\Iop(\tau'){:}$,
\begin{equation}
\int\!\dd\tau\,\hat a^\dagger\hat\Sigma'\hat a
=-\int\!\dd\Omega\,\Omega|\zeta|^2\iint\!\dd\tau\,\dd\tau'\cos[\Omega(\tau'-\tau)]\,{:}\Iop(\tau)\Iop(\tau'){:}
=-\int\!\dd\Omega\,\Omega|\zeta|^2\,{:}\Nop^\dagger\Nop{:}\;;
\end{equation}
(iii) zero, because $V'$ is odd and ${:}\Iop\Iop{:}$ is symmetric. With
$\Nop^\dagger\Nop={:}\Nop^\dagger\Nop{:}+\Nop(0)$ from Eq.~\eqref{eq:L1}, the
mean change for a vacuum THz input is
\begin{equation}
\langle\hat W_{\rm out}\rangle-\langle\hat W_{\rm in}\rangle
=-\hbar\int_0^\infty\!\!\dd\Omega\;\Omega\,|\zeta(\Omega)|^2\,
\langle\Nop^\dagger(\Omega)\Nop(\Omega)\rangle ,
\end{equation}
which is minus the energy of the emitted THz field according to (TN\,10).
The pump pays for the THz radiation by a red shift, with its photon number
unchanged. The part $-\hbar S'(0)\langle\Nop(0)\rangle$, linear in the photon
number, came from the commutator of the THz field with itself: it is the
energy of the spontaneously emitted THz photons, the rectified shot noise of
the theory note.

\subsection{One pump photon}

For a single pump photon with wave function $\psi(\tau)$ and a vacuum THz
input, $\hat D$ acts on each component $|\tau\rangle$ as a displacement of the
THz field and $\ee^{\ii\hat\Phi}$ gives the constant phase $\ee^{-\ii V(0)}$:
\begin{equation}
\hat U|\psi\rangle|0\rangle=\ee^{-\ii V(0)}\int\!\dd\tau\,\psi(\tau)\,|\tau\rangle\,|\beta_\tau\rangle,
\qquad\beta_\tau(\Omega)=-\ii\zeta(\Omega)\,\ee^{\ii\Omega\tau}.
\end{equation}
The number of THz photons is Poisson distributed with mean $B_T$ whatever
$\psi$ is. The mean THz field, $-\ii\zeta\int\dd\tau\,|\psi|^2\ee^{\ii\Omega\tau}$,
is appreciable only at frequencies below the inverse duration of the photon.
A photon much longer than the THz period radiates no mean field but the same
$B_T$ photons: this is spontaneous emission, and it is contained in
(TN\,7)--(TN\,10) without any separate treatment.

\subsection{Numerical tests}

The tests are in the repository under \texttt{code/tests}, in the file for
optical rectification. They use a ladder of $D$ optical frequency bins with
$\Nop=\sum_j\hat a_j^\dagger\hat a_{j+1}$, one THz mode, two pump photons, and
full matrix exponentials. In a finite ladder
$[\Nop,\Nop^\dagger]=\hat n_1-\hat n_D$, so (L4) holds only while the end bins
are empty; the residuals below track their population.

\begin{center}
\begin{tabular}{R{0.24\textwidth}R{0.38\textwidth}R{0.30\textwidth}}
\hline\hline
Equations & Test & Result\\
\hline
(TN\,7)--(TN\,9), with $\Delta kL=2.4$ & state $\ee^{\ii\hat\Phi}\hat D|\psi\rangle$ against the exact $z$-ordered evolution &
 $2\times10^{-5}$ in norm for $D=15$; $0.30$ without $\hat\Phi$\\[4pt]
(TN\,10), (TN\,11) & $\langle\hat b\rangle$, $\langle\hat b^\dagger\hat b\rangle$, and $\langle\Nop\rangle$ at $|\zeta|=0.3$ and $0.6$ &
 better than $10^{-6}$\\[4pt]
(TN\,12)--(TN\,14), with $\Delta k=0$ & pump coherence $\langle\hat a_j^\dagger\hat a_k\rangle$ at $|\zeta|=0.3$ and $0.6$ &
 $7\times10^{-10}$ and $2\times10^{-7}$; $0.03$ and $0.09$ without $\hat\Sigma$\\[4pt]
Energy balance & shift of the mean pump frequency against the THz photon number &
 better than $10^{-9}$\\
\hline\hline
\end{tabular}
\end{center}

Not tested numerically: the pump relation with a phase mismatch, that is,
$\hat\theta$ in Eq.~\eqref{eq:varthetafull}, other than through the state in
the first row; integrals over the THz frequency, since the model has one THz
mode.

\section{Where each assumption enters}

\begin{center}
\begin{tabular}{R{0.36\textwidth}R{0.57\textwidth}}
\hline\hline
Assumption & Used in\\
\hline
$\kappa$ and $\Delta k$ independent of the optical frequency & the generator contains the pump only through $\Nop(\Omega)$, Eq.~\eqref{eq:G}\\[4pt]
no pump dispersion & $\Iop$ depends on $z$ and $t$ only through $\tau$; needed for the same statement\\[4pt]
flux commutes with itself, (L1) & (L4) and (L5), hence every step of Secs.~\ref{sec:U} to \ref{sec:a}\\[4pt]
unbounded band of flat coupling & (L4) in the frequency domain; a band edge gives $[\Nop,\Nop^\dagger]\neq0$\\[4pt]
no pump absorption & unitarity; conservation of $\Iop$\\[4pt]
no coupling among THz modes & $\hat G$ is linear in $\hat b$ and $\hat b^\dagger$, so $[\hat F,\hat F]$ is a number, (L8)\\[4pt]
ordering of the evolution in $z$ & the kernel $V$ of $\hat\Phi$ only; see Sec.~VI\,B of the derivation notes\\
\hline\hline
\end{tabular}
\end{center}

The results (TN\,10) and (TN\,11) do not depend on the last line. The form of
(TN\,12) with a flux-dependent phase and a noise operator is the one known
for self- and cross-phase modulation with a response of finite duration
\cite{boivin1994,shapirorazavi2007}; what is specific here is that the medium
is the THz field and that Eq.~\eqref{eq:half} identifies the phase with the
field at the input and the output of the crystal.

\begin{thebibliography}{9}
\bibitem{blanes2009} S. Blanes, F. Casas, J. A. Oteo, and J. Ros, The Magnus
expansion and some of its applications, Phys. Rep. \textbf{470}, 151 (2009)
[entered from memory, to be checked].
\bibitem{boivin1994} L. Boivin, F. X. K\"artner, and H. A. Haus, Analytical
solution to the quantum field theory of self-phase modulation with a finite
response time, Phys. Rev. Lett. \textbf{73}, 240 (1994) [abstract read; full
text not available; author list from memory].
\bibitem{shapirorazavi2007} J. H. Shapiro and M. Razavi, Continuous-time
cross-phase modulation and quantum computation, New J. Phys. \textbf{9}, 16
(2007) [abstract read].
\end{thebibliography}

\end{document}
